% ==============================================================================
% capitulos/01-introduccion.tex - Introducción del Proyecto
% ==============================================================================

\chapter{Introducción}
\begin{flushleft}
\justifying
El Instituto Politécnico Nacional (IPN), institución rectora de la educación tecnológica en México, establece en su Ley Orgánica \cite{1} y su Reglamento Interno \cite{2} el compromiso de garantizar un entorno de respeto, equidad y legalidad para su comunidad. Ante la identificación de vulneraciones a los derechos de estudiantes, docentes y personal administrativo, el Instituto formalizó la creación de la Defensoría de los Derechos Politécnicos (DDP) mediante su acuerdo de creación oficial \cite{3} . Este órgano, dotado de independencia funcional, constituye el mecanismo institucional responsable de la gestión y seguimiento de inconformidades derivadas de actos u omisiones que contravengan la normativa politécnica.

La administración de estas quejas es un proceso crítico que exige, además de sensibilidad humana, una gestión operativa de alta precisión. Actualmente, la DDP rige su funcionamiento a través del Manual de Organización \cite{4} y el Manual de Procedimientos \cite{5}, instrumentos que definen las etapas de recepción, radicación, investigación y resolución. Sin embargo, la ejecución de estos procesos enfrenta desafíos tecnológicos sustanciales; el flujo de trabajo actual depende de mecanismos manuales y herramientas de oficina descentralizadas que imposibilitan la integración sistémica de la información. Esta carencia de infraestructura digital, sumada al incremento en el volumen de solicitudes, compromete la eficiencia del personal operativo y prolonga los tiempos de respuesta institucionales.

\section{Planteamiento del Problema}
\end{flushleft}
\justifying
La DDP es el órgano institucional encargado de salvaguardar la legalidad y el respeto a los derechos politécnicos de los estudiantes, docentes y personal administrativo del IPN. Su función es fundamental para la atención y resolución de inconformidades; no obstante, ante la creciente demanda de servicios y la complejidad de los procesos operativos actuales, se ha identificado la necesidad de fortalecer su eficiencia mediante la modernización tecnológica. Bajo el flujo de trabajo vigente, se han detectado tres puntos críticos que fundamentan el desarrollo de la solución propuesta en este proyecto.

\subsection{Contexto y situación problemática}
En primer término, se identifica una vulnerabilidad crítica en la estructura operativa de la Defensoría: la desproporción entre la demanda de atención ciudadana y la capacidad del capital humano especializado. Actualmente, la función sustantiva de protección legal y resolución de controversias recae de forma exclusiva sobre un cuerpo técnico de apenas cinco abogados, quienes tienen la responsabilidad de gestionar la totalidad de las quejas declaradas procedentes dentro de una comunidad que abarca cientos de miles de integrantes.

Bajo el modelo de gestión vigente, este personal de alta especialización jurídica se ve absorbido por una densidad administrativa abrumadora. Al carecer de una plataforma centralizada, los abogados deben dedicar una fracción significativa de su jornada a tareas procedimentales, tales como el registro manual de expedientes, el seguimiento analógico de plazos procesales y la organización descentralizada de evidencias. Esta carga de trabajo no sustantiva actúa como un factor de erosión del tiempo efectivo, limitando el análisis profundo que cada caso amerita para garantizar el debido proceso.

En consecuencia, la implementación de una solución tecnológica avanzada no debe entenderse simplemente como una herramienta de apoyo, sino como un imperativo estratégico de modernización. Al automatizar los flujos de trabajo repetitivos y digitalizar la trazabilidad de los casos, se faculta a los juristas para recuperar su capacidad de enfoque en la resolución sustantiva. El objetivo es transitar de un modelo de gestión reactivo y manual hacia un ecosistema digital eficiente, donde el talento humano sea liberado de la burocracia técnica para centrarse íntegramente en la salvaguarda de los derechos de la comunidad politécnica.

\subsection{Problema a resolver}
Derivado de lo anterior, existe un reto en la eficiencia de la recepción y clasificación de solicitudes durante la etapa de \textit{``Primer Contacto''}. Ya que, al recibir información por múltiples canales descentralizados, las solicitudes suelen presentarse de forma incompleta o no son competentes para el órgano de la DDP, resaltando la falta de un asistente tecnológico para el filtrado de competencia, obligando a la revisión manual, retrasando la atención a casos de alta prioridad.

Finalmente, la gestión de documentos realizada de forma tradicional y la falta de seguimiento impactan en la transparencia del proceso. La actual dependencia de métodos de archivo físicos dificulta la consulta rápida de antecedentes. La transición a una plataforma integral no solo brindaría certidumbre al quejoso sobre el estatus de su queja, sino que permitiría al personal de la DDP la revisión de documentos y datos importantes sobre los casos.

\section{Solución Propuesta}
Se propone el desarrollo de una plataforma web integral diseñada para centralizar, automatizar y optimizar el ciclo de vida de las quejas dentro de la DDP. La solución se fundamenta en una arquitectura basada en microservicios y contempla la implementación de un módulo de asistencia inteligente en la etapa de primer contacto, el cual utiliza técnicas de PLN para asistir al usuario en la determinación preliminar de la competencia de su queja. Se permitirá la gestión basada en roles, garantizando que cada actor interactúe con la información de acuerdo con sus facultades normativas, asegurando el seguimiento total del expediente desde su recepción hasta su resolución final.

\section{Justificación}
La implementación de esta plataforma web se justifica por la necesidad estratégica de fortalecer los procesos operativos de la DDP, transformando una gestión actual basada en métodos manuales y herramientas de oficina descentralizadas hacia un flujo digital inteligente.

\subsection{Ámbito social e institucional}
En el ámbito social e institucional, el proyecto busca fortalecer varios puntos empezando por la eliminación de barreras que impiden una mejor comunicación entre la comunidad politécnica y la DDP al proporcionar un canal de seguimiento transparente que elimine la incertidumbre del quejoso. De igual forma, al automatizar las tareas administrativas, se potencia la capacidad de respuesta del personal de la Defensoría y se reduce el costo de gestión operativa, permitiendo que el cuerpo de abogados centralice sus esfuerzos en el análisis jurídico de fondo a las quejas procedentes, apoyando una atención equitativa y eficiente ante la creciente demanda de quejas que reciben diariamente.

\subsection{Perspectiva técnica}
Desde la perspectiva técnica, el desarrollo de este sistema se fundamenta en la adopción de una arquitectura de microservicios permitiendo la oportunidad de escalabilidad y mantenimiento independiente de sus componentes. Ya que la integración de tecnologías como Java Spring Boot y PostgreSQL permite el manejo seguro de información sensible, mientras que la incorporación de un clasificador basado en PLN demuestra una innovación clave para resolver los puntos significativos desde el área de primer contacto. Esta solución resuelve la fragmentación de datos y fortalece la integridad de un repositorio histórico de quejas para futuras consultas y auditorias.

\section{Objetivos}

\subsection{Objetivo General}
Desarrollar una plataforma web basada en una arquitectura de microservicios mediante la integración de módulos de gestión documental y un motor de procesamiento de lenguaje natural para sistematizar el flujo operativo de recepción, clasificación y seguimiento de quejas de la DDP.

\subsection{Objetivos específicos}
\begin{itemize}
    \item Modelar las etapas de proceso y resolución de quejas con base en las reglas de negocio del Manual de Procedimientos de la DDP para automatizar la transición de estados del expediente y el cumplimiento de términos legales.
    \item Diseñar una arquitectura de microservicios desacoplada empleando el ecosistema de Java Spring Boot y comunicación vía API REST para garantizar la independencia operativa entre el sistema de gestión administrativa y el módulo de inteligencia artificial.
    \item Implementar un clasificador de texto neuro-simbólico utilizando Word Embeddings y sistemas basados en reglas para determinar la competencia de las solicitudes desde el primer contacto y reducir la carga de clasificación manual.
    \item Construir un repositorio digital de evidencias y expedientes mediante el sistema de gestión de bases de datos PostgreSQL para centralizar la documentación soporte y asegurar la integridad de la información jurídica manejada por la Defensoría.
    \item Desarrollar un módulo de consulta y seguimiento para el quejoso a través de un sistema de folios digitales y notificaciones de estatus para brindar certidumbre sobre el avance del proceso y reducir la saturación de los canales de comunicación presenciales.
\end{itemize}

\section{Alcances y Limitaciones}

\subsection{Alcances}
La plataforma web permitirá el registro, seguimiento y gestión integral de las quejas recibidas por la DDP, abarcando desde el primer contacto hasta la resolución final del expediente. El sistema implementará un clasificador basado en PLN para asistir en la determinación de competencia de las solicitudes, así como un repositorio digital para el almacenamiento seguro de documentos y evidencias. Se desarrollarán módulos específicos para cada rol institucional (Abogado Asesor, Subdefensor, Recepcionista, Área de Primer Contacto y Titular) y se proporcionará un sistema de consulta pública mediante folios digitales.

\subsection{Limitaciones}
El sistema no sustituirá el criterio jurídico de los abogados de la DDP, sino que fungirá como una herramienta de apoyo en la clasificación inicial. La plataforma no realizará resoluciones automáticas de quejas ni emitirá dictámenes legales. Asimismo, el clasificador basado en IA requerirá un período de entrenamiento y ajuste con datos reales para alcanzar niveles óptimos de precisión, por lo que en etapas iniciales operará de manera asistida con supervisión humana.
